% =============================================================
%  Appendix A3 -- Wild Cluster Bootstrap Details
%  Label: app:bootstrap
% =============================================================

\section{Wild Cluster Bootstrap}\label{app:bootstrap}

\subsection{Methodology}

Asymptotic cluster-robust standard errors can be unreliable
with fewer than approximately 50 clusters
\citep{Cameron2008, MacKinnon2017}.
With 31 province clusters, I supplement conventional inference
with the wild cluster bootstrap.
This appendix describes the methodology in detail and reports
the full results.

The wild cluster bootstrap constructs the null distribution of
the test statistic by resampling residuals at the cluster level.
Two versions of the procedure exist: the wild cluster
unrestricted (WCU) and the wild cluster restricted (WCR).
In the WCU version, the bootstrap samples are constructed
from residuals of the unrestricted model (which includes the
coefficient of interest).
In the WCR version, the bootstrap samples are constructed
from residuals of the restricted model, which imposes the
null hypothesis $H_0\colon \beta = 0$.

I use the WCR version throughout.
\citet{MacKinnon2017} show that the restricted bootstrap
has superior size control when the number of clusters is
small.
The intuition is that the restricted residuals are computed
under the null, so the bootstrap distribution is centred at
the correct location under $H_0$.
The unrestricted version uses residuals computed under the
alternative, which can introduce a bias in the bootstrap
distribution when the number of clusters is small and the
test statistic is far from zero.
With 31 clusters and a $t$-statistic exceeding 5 for the
forecast-error coefficient, this bias is potentially
nontrivial, making the restricted version the more
conservative choice.

I use the six-point weight distribution proposed by
\citet{Webb2023}, which is designed for settings with few
clusters.
The standard wild bootstrap draws weights from a two-point
distribution---either Rademacher $\{-1, +1\}$ or Mammen
$\{-(1{-}\sqrt{5})/2, (1{+}\sqrt{5})/2\}$---with equal
probability.
With $G = 31$ clusters, the Rademacher distribution produces
$2^{31} \approx 2.15 \times 10^9$ distinct bootstrap samples.
While this is a large number, the effective coverage of the
null distribution can be uneven because each bootstrap draw
assigns the same weight to all observations within a cluster.
The six-point distribution provides a richer set of weights
$\{w_1, \ldots, w_6\}$ with $6^{31} \approx 1.3 \times 10^{24}$
distinct samples, which yields a finer approximation to the
true null distribution.
\citet{Webb2023} show in simulations with 10--30 clusters that
the six-point distribution produces rejection rates closer to
the nominal size than either the Rademacher or Mammen
distributions.

Each bootstrap replication proceeds as follows.

\begin{algorithm}[htbp]
\caption{Wild Cluster Restricted Bootstrap (WCR)}\label{alg:wcr_bootstrap}
\begin{algorithmic}[1]
\Require Observed data $\{(y_{ig}, X_{ig})\}$; number of clusters $G = 31$; replications $B = 9{,}999$
\Ensure Bootstrap $p$-value for $H_0\colon \beta = 0$
\State Estimate the restricted model (imposing $H_0\colon \beta = 0$) and obtain residuals $\hat{\varepsilon}_{ig}^{r}$ and fitted values $\hat{y}_{ig}^{r}$ for each observation $i$ in cluster $g$
\State Estimate the unrestricted model on the original data and record $t_{\text{obs}}$
\For{$b = 1, \ldots, B$}
    \For{$g = 1, \ldots, G$}
        \State Draw $w_g$ independently from the six-point distribution (each value with probability $1/6$)
    \EndFor
    \For{each observation $i$ in cluster $g$}
        \State $\varepsilon^*_{ig} \gets w_g \cdot \hat{\varepsilon}_{ig}^{r}$ \Comment{Preserve within-cluster correlation}
        \State $y^*_{ig} \gets \hat{y}_{ig}^{r} + \varepsilon^*_{ig}$
    \EndFor
    \State Re-estimate the unrestricted model on $(y^*_{ig}, X_{ig})$
    \State Record $t^*_b$ for the coefficient of interest
\EndFor
\State \Return $\hat{p} = \dfrac{1}{B+1}\displaystyle\sum_{b=1}^{B} \mathbf{1}\!\left(|t^*_b| \geq |t_{\text{obs}}|\right)$
\end{algorithmic}
\end{algorithm}

The choice of $B = 9{,}999$ is standard in the literature and
ensures that the bootstrap $p$-value has a precision of
approximately $\pm 0.01$ at the 5 percent level.
With $B = 9{,}999$ and a true rejection rate of 0.05, the
Monte Carlo standard error of the estimated $p$-value is
$\sqrt{0.05 \times 0.95 / 10{,}000} \approx 0.002$.

\subsection{Results}

Table~\ref{tab:bootstrap_results} reports the bootstrap
results for the two main coefficients (expectation and
forecast error) under both fixed-effect structures.

\begin{table}[htbp]
\centering
\caption{Wild Cluster Bootstrap Results}
\label{tab:bootstrap_results}
\begin{threeparttable}
\small
\setlength{\tabcolsep}{4pt}
\begin{tabular}{llccccc}
\toprule
Equation & Fixed effects & $\hat{\beta}$ & SE &
  Asymptotic $p$ & Bootstrap $p$ & $N$ \\
\midrule
Expectations & Province + wave & 0.090 & 0.053
  & 0.100 & 0.088 & 90,133 \\
Expectations & Region $\times$ wave & 0.242 & 0.174
  & 0.174 & 0.237 & 90,133 \\
Expectations & Province + region $\times$ wave & 0.119 & 0.050
  & 0.024 & 0.024 & 90,133 \\[4pt]
Forecast error & Province + wave & $-5.351$ & 1.201
  & $<0.001$ & $<0.001$ & 69,533 \\
Forecast error & Region $\times$ wave & $-5.555$ & 1.198
  & $<0.001$ & $<0.001$ & 69,533 \\
Forecast error & Province + region $\times$ wave & $-7.127$ & 1.193
  & $<0.001$ & $<0.001$ & 69,533 \\
\bottomrule
\end{tabular}
\begin{tablenotes}[flushleft]
\footnotesize
\item Webb six-point distribution with 9,999 replications.
  Restricted bootstrap (WCR) imposing $H_0\colon \beta = 0$.
\item $p$-values are two-sided.
\end{tablenotes}
\end{threeparttable}
\end{table}

The forecast-error coefficient $\hat{\beta}_3$ is robust to
the bootstrap correction under both fixed-effect structures.
The bootstrap $p$-values for Eq~3 are below $0.001$ in both
cases, matching the asymptotic $p$-values.
The 95-percent bootstrap confidence intervals for
$\hat{\beta}_3$ exclude zero by a wide margin under both
specifications.

The expectation coefficient $\hat{\beta}_1$ is marginal.
Under province-plus-wave FE, the bootstrap $p$-value of
$0.088$ is slightly below the conventional 10 percent
threshold, while the asymptotic $p$-value of $0.100$ is at
the boundary.
Under region-by-wave FE, both the asymptotic ($0.174$) and
bootstrap ($0.237$) $p$-values are above 10 percent.
The bootstrap does not reverse any asymptotic conclusion but
does widen the $p$-value modestly for Eq~1 under
region-by-wave FE, reflecting the conservative nature of the
WCR procedure with few clusters.

\subsection{Power Considerations}\label{app:bootstrap_power}

The wild cluster bootstrap corrects the size of the test
(the probability of rejecting a true null) but does not
improve its power (the probability of rejecting a false null).
With 31 clusters, the effective degrees of freedom for
cluster-level variation are limited, and power against small
to moderate effect sizes is low.

This power limitation is relevant for interpreting the
imprecise expectation coefficient.
The identifying variation in $\hat{\beta}_1$ operates at the
province-wave level: with 31 provinces and 4 waves, the
maximum number of province-wave cells is 124, but the
effective number is reduced by the fixed-effect structure.
Under region-by-wave FE with approximately 24 region-wave
cells, the residual degrees of freedom for the meat-shock
coefficient are approximately $124 - 24 = 100$ at the
cell level, but cluster-robust inference treats the 31
provinces as the unit, yielding an effective sample size of
31 for the purpose of the $t$-distribution.
Power calculations for a two-sided test at the 5 percent
level with 30 degrees of freedom suggest that detecting an
effect of the size estimated for $\hat{\beta}_1$ (0.242
with s.e. 0.174, implying a $t$-ratio of 1.4) requires
approximately 50--60 clusters to achieve 80 percent power.
The current design with 31 clusters is underpowered for the
expectation coefficient.

The forecast-error coefficient does not face this power
problem.
The $t$-ratio for $\hat{\beta}_3$ exceeds 5 under both
FE structures, which is well above any critical value for
30 degrees of freedom.
The bootstrap confirms that the apparent power is not an
artifact of size distortion: the $p$-values remain below
$0.001$ after the correction.

The asymmetry in power between Eq~1 and Eq~3 reflects both
the measurement properties of the dependent variables and the
structure of the identifying variation.
The ordinal expectation in Eq~1 takes three values
($-1, 0, 1$), which limits the effective variation in the
dependent variable.
The forecast error in Eq~3 combines the ordinal expectation
with a continuous CPI measure, producing a dependent variable
with substantially more dispersion.
The greater dispersion in the dependent variable of Eq~3
translates to a higher signal-to-noise ratio for the
coefficient of interest and, consequently, a larger
$t$-statistic for a given effect size.

A further consideration is that the wild cluster bootstrap
is designed for inference, not for estimation.
The point estimates $\hat{\beta}_1$ and $\hat{\beta}_3$ are
OLS estimates and are unbiased under the standard assumptions.
The bootstrap corrects only the distribution of the test
statistic.
The fact that the bootstrap widens the $p$-value for
$\hat{\beta}_1$ under region-by-wave FE (from 0.174 to 0.237)
does not imply that the point estimate is biased; it implies
that the asymptotic standard error slightly understates the
true sampling variability of the $t$-ratio with 31 clusters.

These power considerations motivate the emphasis on the
forecast-error equation as the primary test.
The forecast-error coefficient provides the most powerful
test of overreaction because it combines the expectation
and CPI channels into a single dependent variable with
high dispersion, and because the sign of $\hat{\beta}_3$
alone is sufficient to reject the rational benchmark without
requiring a precise estimate of the expectation response.
